% ---- score2 standard preamble (from: python3 score2_tools.py preamble) ----
\documentclass[10pt]{extarticle}
\usepackage[a4paper,margin=1.5cm]{geometry}
\usepackage[T1]{fontenc}
\usepackage{amsmath,amssymb}
\usepackage{graphicx}
\usepackage[export]{adjustbox}
\usepackage{multicol}
\usepackage{enumitem}
\usepackage{booktabs,array,tabularx}
\setlength{\columnsep}{0.7cm}
\setlength{\parindent}{0pt}
\setlength{\parskip}{3pt}
\pagestyle{empty}
\makeatletter
% figure registry (figdefs.tex is generated by the extract step)
\newcommand{\qfigdef}[4]{\@namedef{qf@file@#1}{#2}\@namedef{qf@w@#1}{#3}\@namedef{qf@wide@#1}{#4}}
\newcommand{\qf@check}[1]{\@ifundefined{qf@file@#1}{\PackageError{score2}{Unknown figure id #1}{}}{}}
\newcommand{\qf@img}[4][]{\includegraphics[width=#2\dimexpr\@nameuse{qf@w@#4}pt\relax,max width=#3,max height=0.45\textheight,#1]{\@nameuse{qf@file@#4}}}
% Full-width questions: finalize wraps every question that contains a wide figure in \qfullbegin...\qfullend,
% so the WHOLE question (stem, figures, options) is set across the page in reading order, then the two columns
% resume. The question's head - its number, stem and first figure (\qfullsplit marks the end of the head) -
% is kept on one page, so a number is never left at a page foot with its figure overleaf; the rest may break
% across pages normally, so no big gaps are left. Without a \qfullsplit, the whole question is kept together
% if it is under 45% of a page. Never write these yourself.
\newif\ifqf@full
\newif\ifqf@boxed
\newsavebox{\qf@box}
\newcommand{\qf@open}{\global\qf@fulltrue\par\global\qf@boxedtrue
  \setbox\qf@box\vbox\bgroup\hsize\textwidth\linewidth\textwidth}
\newcommand{\qf@emit}[1]{\par\egroup\global\qf@boxedfalse
  \ifdim\dimexpr\ht\qf@box+\dp\qf@box\relax<#1\textheight\noindent\usebox{\qf@box}\par\else\unvbox\qf@box\fi}
\newcommand{\qfullsplit}{\ifqf@boxed\qf@emit{0.85}\fi}
\newcommand{\qf@close}{\ifqf@boxed\qf@emit{0.45}\fi\global\qf@fullfalse}
\newcommand{\qfullbegin}{\end{multicols}\qf@open}
\newcommand{\qfullend}{\qf@close\begin{multicols}{2}}
\newcommand{\qfullstart}{\qf@open}
\newcommand{\qfulldone}{\qf@close}
% \qfig[scale]{id}: block figure, centred, directly where it is written (size set by the tool).
% Fallback if a wide figure is not inside a full-width question: break out of the columns in place - no floats.
\newcommand{\qfig}[2][1]{\qf@check{#2}%
  \ifqf@full\qf@blk{#1}{#2}\else
  \ifnum\@nameuse{qf@wide@#2}=1 \expandafter\@firstoftwo\else\expandafter\@secondoftwo\fi
  {\end{multicols}\begin{center}\qf@img{#1}{\textwidth}{#2}\end{center}\begin{multicols}{2}}%
  {\qf@blk{#1}{#2}}\fi}
\newcommand{\qf@blk}[2]{\par\smallskip{\centering\qf@img{#1}{\linewidth}{#2}\par}\smallskip}
% \qfigi[scale]{id}: inline figure (inside option lists / tables / minipages); never breaks columns.
\newcommand{\qfigi}[2][1]{\qf@check{#2}\qf@img[valign=c]{#1}{\linewidth}{#2}}
\makeatother
% structure macros - lint reads these to check completeness and figure placement
\newcommand{\papersec}[1]{\par\bigskip{\centering\Large\bfseries PAPER-#1\par}\medskip}
\newcommand{\partsec}[2]{\par\medskip\noindent{\large\bfseries PART-#1: #2}\par\smallskip}
\newcommand{\sect}[1]{\par\smallskip\noindent\textbf{#1}\par}
\newcommand{\qn}[1]{\par\smallskip\noindent\textbf{#1.}~}
\newcommand{\paraq}[2]{\par\smallskip\noindent\textbf{Paragraph for Questions #1 and #2}\par}
\newenvironment{qoptions}{\begin{itemize}[leftmargin=1.8em,itemsep=1pt,topsep=2pt,label={}]}{\end{itemize}}
\newcommand{\opt}[2]{\item[(#1)] #2}
% figure registry (generated by score2_tools.py finalize)
\qfigdef{p1_phy_s1_q03_1}{Score-I - Test 10 (24 Feb 2022)-img/p1_phy_s1_q03_1.png}{133.5}{0}
\qfigdef{p1_phy_s1_q06_1}{Score-I - Test 10 (24 Feb 2022)-img/p1_phy_s1_q06_1.png}{102.1}{0}
\qfigdef{p1_phy_s1_q08_1}{Score-I - Test 10 (24 Feb 2022)-img/p1_phy_s1_q08_1.png}{160.1}{0}
\qfigdef{p1_phy_s2_q01_1}{Score-I - Test 10 (24 Feb 2022)-img/p1_phy_s2_q01_1.png}{120.5}{0}
\qfigdef{p1_phy_s2_q02_1}{Score-I - Test 10 (24 Feb 2022)-img/p1_phy_s2_q02_1.png}{131.3}{0}
\qfigdef{p1_phy_s2_q03_1}{Score-I - Test 10 (24 Feb 2022)-img/p1_phy_s2_q03_1.png}{128.3}{0}
\qfigdef{p1_phy_s2_q06_1}{Score-I - Test 10 (24 Feb 2022)-img/p1_phy_s2_q06_1.png}{171.1}{0}
\qfigdef{p1_phy_s3_q02_1}{Score-I - Test 10 (24 Feb 2022)-img/p1_phy_s3_q02_1.png}{70.8}{0}
\qfigdef{p1_chem_s1_q01_A}{Score-I - Test 10 (24 Feb 2022)-img/p1_chem_s1_q01_A.png}{77.7}{0}
\qfigdef{p1_chem_s1_q01_B}{Score-I - Test 10 (24 Feb 2022)-img/p1_chem_s1_q01_B.png}{79.8}{0}
\qfigdef{p1_chem_s1_q01_C}{Score-I - Test 10 (24 Feb 2022)-img/p1_chem_s1_q01_C.png}{81.5}{0}
\qfigdef{p1_chem_s1_q01_D}{Score-I - Test 10 (24 Feb 2022)-img/p1_chem_s1_q01_D.png}{82.8}{0}
\qfigdef{p1_chem_s1_q04_1}{Score-I - Test 10 (24 Feb 2022)-img/p1_chem_s1_q04_1.png}{204.0}{0}
\qfigdef{p1_chem_s1_q04_2}{Score-I - Test 10 (24 Feb 2022)-img/p1_chem_s1_q04_2.png}{52.7}{0}
\qfigdef{p1_chem_s1_q04_D}{Score-I - Test 10 (24 Feb 2022)-img/p1_chem_s1_q04_D.png}{52.9}{0}
\qfigdef{p1_chem_s1_q05_A}{Score-I - Test 10 (24 Feb 2022)-img/p1_chem_s1_q05_A.png}{53.6}{0}
\qfigdef{p1_chem_s1_q05_B}{Score-I - Test 10 (24 Feb 2022)-img/p1_chem_s1_q05_B.png}{30.8}{0}
\qfigdef{p1_chem_s1_q05_D}{Score-I - Test 10 (24 Feb 2022)-img/p1_chem_s1_q05_D.png}{35.6}{0}
\qfigdef{p1_chem_s2_q03_1}{Score-I - Test 10 (24 Feb 2022)-img/p1_chem_s2_q03_1.png}{103.5}{0}
\qfigdef{p1_chem_s3_pre_A}{Score-I - Test 10 (24 Feb 2022)-img/p1_chem_s3_pre_A.png}{127.6}{0}
\qfigdef{p1_chem_s3_pre_B}{Score-I - Test 10 (24 Feb 2022)-img/p1_chem_s3_pre_B.png}{156.2}{0}
\qfigdef{p1_chem_s3_pre_C}{Score-I - Test 10 (24 Feb 2022)-img/p1_chem_s3_pre_C.png}{102.3}{0}
\qfigdef{p1_chem_s3_pre_D}{Score-I - Test 10 (24 Feb 2022)-img/p1_chem_s3_pre_D.png}{114.1}{0}
% ---- body skeleton ----
% \begin{document}
% \begin{center}{\Large\bfseries JEE (Main+Advanced) Score-II --- Test NN (DD Mon YYYY)}\end{center}
% \begin{multicols}{2}
%   \papersec{1}   (only for two-paper tests)
%   \partsec{1}{PHYSICS} \sect{Section-I (i): Single Correct} \qn{1} ... \qfig{p1_phy_q01_1} ...
% \end{multicols}
% \end{document}

\begin{document}
\raggedcolumns
\setlength{\tabcolsep}{3pt}
\begin{center}{\Large\bfseries JEE (Main+Advanced) Score-I --- Test 10 (24 Feb 2022)}\end{center}
\begin{multicols}{2}
\raggedcolumns
\partsec{1}{PHYSICS}
\sect{SECTION-I}
\qn{1} If a black metal foil is warmed by radiation by a sphere at temperature $T$ and black metal foil is kept at a distance $d$ from the sphere and power received by black metal foil is $P$ :-
\begin{qoptions}
\opt{A}{If temperature of sphere is increased by $T$ then power received by metal is $2P$}
\opt{B}{If temperature of sphere is increased by $T$ then power received by metal is $16P$}
\opt{C}{If temperature of sphere becomes $2T$ and distance is $2d$ then power received by metal is $4P$}
\opt{D}{All are correct}
\end{qoptions}

\qn{2} Consider a region where charge density varies with radius as:- (here $\rho_0$ is a constant)\\
$\rho=\rho_0 r\quad 0\le r\le R$\\
$\phantom{\rho}=0\quad r>R$\\
Choose the correct statement (s) :-
\begin{qoptions}
\opt{A}{Electric field varies as $r^2$ for $r<R$}
\opt{B}{Electric field varies as $\dfrac1{r^2}$ for $r>R$}
\opt{C}{Electric field varies as $r$ for $r<R$}
\opt{D}{Electric field varies as $\dfrac1r$ for $r>R$}
\end{qoptions}

\qn{3} Four identical capacitor each of capacitance 5 $\mu$F are connected as shown in Figure. They are attached to the battery of emf $E=9$ V. Initially the key $K_2$ is open, and the switch $K_1$ is closed. Then key $K_1$ is opened and $K_2$ is closed.
\qfig{p1_phy_s1_q03_1}
\begin{qoptions}
\opt{A}{The initial potential difference across capacitor 2 is 3 V.}
\opt{B}{The final potential difference across capacitor 2 is 1.8 V.}
\opt{C}{The charge flown through the battery after the reversal of switches is 3 $\mu$C.}
\opt{D}{The charge flown through the battery after the reversal of switches is zero.}
\end{qoptions}

\qn{4} A hydrogen like atom of atomic number $Z$ is in excited state of quantum number $2n$. It can emit a maximum energy photon of 204 eV. If it makes a transition to a quantum number $n$, a photon of energy 40.8 eV is emitted.
\begin{qoptions}
\opt{A}{The value of $n$ is 2.}
\opt{B}{The value of $Z$ is 3.}
\opt{C}{The minimum energy that can be emitted by this atom during de-exitation is 10.57 eV.}
\opt{D}{The ground state energy for this atom is $-217.6$ eV.}
\end{qoptions}

\qn{5} Mark the INCORRECT statement(s) for a standing wave established in a string :
\begin{qoptions}
\opt{A}{Energy of an element of medium remains constant.}
\opt{B}{Phase difference between any two particles can be $\dfrac\pi2$}
\opt{C}{Power can get transferred through antinodes.}
\opt{D}{Power can get transferred through nodes.}
\end{qoptions}

\qn{6} One mole of an ideal gas undergoes a thermodynamic cyclic process as shown in the figure. The cyclic process consists of an isochoric, an adiabatic process and an isothermal (may not be in the same sequence).
\qfig{p1_phy_s1_q06_1}
\begin{qoptions}
\opt{A}{The ratio of molar specific heat at constant pressure to that at constant volume is $\dfrac{\ell n3}{\ell n3-\ell n2}$.}
\opt{B}{The ratio of molar specific heat at constant pressure to that at constant volume is $\dfrac{\ell n2}{\ell n3-\ell n2}$}
\opt{C}{Magnitude of work done by the gas in the isothermal process is $2P_0V_0\,\ell n(3/2)$.}
\opt{D}{Magnitude of work done by the gas in the isothermal process is $3P_0V_0\,\ell n(3/2)$.}
\end{qoptions}

\qn{7} A student uses a simple pendulum of exactly 1 m length to determine $g$, the acceleration due to gravity. He uses a stop watch with the least count of 1 sec for this and records 40 seconds for 20 oscillations. For this observation, which of the following statement(s) is (are) true?
\begin{qoptions}
\opt{A}{Error $\Delta T$ in measuring $T$, the time period is 0.05 seconds}
\opt{B}{Error $\Delta T$ in measuring $T$, the time period is 1 second}
\opt{C}{Percentage error in the determination of $g$ is 5\%}
\opt{D}{Percentage error in the determination of $g$ is 2.5\%}
\end{qoptions}

\qn{8} The refractive index of the medium with a certain region, $x>0$, $y>0$, changes with $y$. A thin light ray travelling in the $x$-direction in medium having refractive index $\mu_0=1$ strikes another medium of refractive index $\mu$ at right angles and moves through the medium along a circular arc of radius $R$ as shown in the figure. The material with the greatest known refractive index is diamond, but even the refractive index of this material does not reach the value $\mu_{\max}=2.5$. It is this limit that sets the maximum angular size of the arc the light ray can cover. Angular size of arc is the angle subtended by the arc at the centre. Which of the following statement(s) is/are \textbf{CORRECT} ?
\qfig{p1_phy_s1_q08_1}
\begin{qoptions}
\opt{A}{The variation of refractive index $\mu$ with $y$ is given as $\mu=\dfrac{R}{R-y}$.}
\opt{B}{The unit vector in the direction of refracted light at $y=\dfrac R2$ is $\dfrac12\hat i+\dfrac{\sqrt3}{2}\hat j$.}
\opt{C}{The maximum angular size of the arc of light is $\theta_{\max}$. The value of $\cos\theta_{\max}=2/5$}
\opt{D}{The maximum angular size of the arc of light is $\theta_{\max}$. The value of $\sin\theta_{\max}=2/5$}
\end{qoptions}

\sect{SECTION-II}
\qn{1} The ratio $\left(\dfrac{m_1}{m_2}\right)$ for which the massless rod remains horizontal is :-
\qfig{p1_phy_s2_q01_1}

\qn{2} A solid cylinder is kept over a rough plank which is oscillating along $x$-axis according to equation $x=A\cos(10)t$, where $x$ is in m and $t$ is in s. If coefficient friction between cylinder and plank is 0.3, then find maximum amplitude of plank (in cm) possible, so that cylinder never slips on plank as it rolls on plank.
\qfig{p1_phy_s2_q02_1}

\qn{3} Consider the system shown in figure, with two equal masses $m=1$ kg and a spring with spring constant 800 N/m. The coefficient of kinetic friction is $\mu=0.2$ and the pulley is frictionless. The system is held with the spring at its relaxed length and then released. Find maximum displacement (in cm) of the block in subsequent motion.
\qfig{p1_phy_s2_q03_1}

\qn{4} A vertical capillary tube with inside radius 0.25 mm is submerged into water so that the length of its part protruding over the water surface is equal to 25 mm. Surface tension of water is $73\times10^{-3}$ N/m and angle of contact is zero degree for glass and water, acceleration due to gravity is 9.8 m/s$^2$. Then value of $10R$ approximately (in mm) is (where $R$ is radius of meniscus and $h$ is height of water in capillary tube)

\qn{5} The r.m.s value of current $i=3+4\sin\left(\omega t+\dfrac\pi3\right)$ is $\left(\sqrt{1+x^2}\right)$A. Write value of $x$.

\qn{6} Circuit shows seven identical bulbs (A to G) connected through a battery of emf 200 V. The bulbs are rated as 200 V, 100 Watt. The power dissipated in the bulb which glow brightest is given by $P$ watt. Fill the value of $\dfrac23\sqrt P$.
\qfig{p1_phy_s2_q06_1}

\qn{7} A small area is removed from a uniform spherical shell of mass $M$ and radius $R$, then the gravitational field intensity near the hollow portion is $GM/nR^2$. Find $n$ :

\qn{8} Cart of mass $m$ moving with a velocity 12 m/s to the right collides perfectly elastically with a cart of mass 4 kg that is originally at rest. After the collision, the cart of mass $m$ moves to the left with a velocity of 6 m/s. Assuming the collision to be head on, what is the velocity (in m/s) of the centre of mass of the two carts before the collision ?

\sect{SECTION-IV}
\qn{1} A charged particle with some initial velocity is projected in a region where uniform electric and/or magnetic fields are present. In Column-I information about the existence of electric and/or magnetic field and direction of initial velocity of charged particle are given, while in column II the possible paths of charged particle is mentioned. Match the entries of Column I with the entries of Column II.
\par\nobreak
\begin{center}\footnotesize
\begin{tabular}{|p{0.5cm}|p{3.5cm}|p{0.6cm}|p{2.4cm}|}
\hline
\multicolumn{2}{|c|}{Column--I} & \multicolumn{2}{c|}{Column--II}\\ \hline
(A) & $\vec E=0,\ \vec B\ne0$ and initial velocity is at an unknown angle with $\vec B$ & (P) & Straight line\\ \hline
(B) & $\vec E\ne0,\ \vec B=0$ and initial velocity is at an unknown angle with $\vec E$ & (Q) & Parabola\\ \hline
(C) & $\vec E\ne0,\ \vec B\ne0,\ \vec E\parallel\vec B$ and initial velocity is perpendicular to $\vec E$ & (R) & Circular\\ \hline
(D) & $\vec E\ne\vec B,\ \vec B\ne0,\ \vec E$ perpendicular $\vec B$ and initial velocity is perpendicular to both $\vec E$ and $\vec B$ & (S) & Helical path with nonuniform pitch\\ \hline
 & & (T) & Helical path with uniform pitch\\ \hline
\end{tabular}
\end{center}

\qn{2} In the standard YDSE apparatus shown in figure ; $\Delta x$ is the path difference between $S_2P$ \& $S_1P$. Now a glass slab is introduced in front of $S_2$. Then match the following :-
\qfig{p1_phy_s3_q02_1}
\par\nobreak
\begin{center}\small
\begin{tabular}{|p{0.5cm}|p{3.2cm}|p{0.6cm}|p{2.4cm}|}
\hline
 & Column-I & & Column-II\\ \hline
(A) & $\Delta x$ at P will & (P) & Increases\\ \hline
(B) & Fringe width will & (Q) & Shift downward\\ \hline
(C) & Fringe pattern will & (R) & Remain same\\ \hline
(D) & Number of fringes between O \& P will & (S) & Shift upward\\ \hline
 & & (T) & Decrease\\ \hline
\end{tabular}
\end{center}

\partsec{2}{CHEMISTRY}
\sect{SECTION-I}
\qn{1} Select the correct curve(s) :\\
If $v$ = velocity of electron in Bohr's orbit\\
$r$ = Radius of electron in Bohr's orbit\\
P.E. = Potential energy of electron in Bohr's orbit\\
K.E. = Kinetic energy of electron in Bohr's orbit.
\begin{qoptions}
\opt{A}{\qfigi{p1_chem_s1_q01_A}}
\opt{B}{\qfigi{p1_chem_s1_q01_B}}
\opt{C}{\qfigi{p1_chem_s1_q01_C}}
\opt{D}{\qfigi{p1_chem_s1_q01_D}}
\end{qoptions}

\qn{2} 200 mL of 0.01 M NaOH is mixed with 400 mL of 0.005 M $\mathrm{H_2S}$ solution ($\mathrm{p}K_1$ and $\mathrm{p}K_2$ of $\mathrm{H_2S}$ are 6 \& 12 respectively)
\begin{qoptions}
\opt{A}{Resulting solution is basic}
\opt{B}{pH of solution is 9}
\opt{C}{pOH of solution is 5}
\opt{D}{pH of solution is nearly 7}
\end{qoptions}

\qn{3} A metal exists in two allotropic forms. $\alpha$-form is BCC and $\beta$-form is FCC. The distance between nearest neighbours is same in both forms. Which of the following is/are \textbf{INCORRECT} information(s) related with these two forms of the metal ?
\begin{qoptions}
\opt{A}{The distances between next nearest neighbours is same in both.}
\opt{B}{Percentage void space in $\alpha$-form is greater than that in $\beta$-form.}
\opt{C}{The diameter of largest sphere which may be fitted in the crystals of metal, without disturbing the distance between nearest neighbours, is same in both forms.}
\opt{D}{The cell constant is higher for the unit cell of $\beta$-form.}
\end{qoptions}

\qn{4} Which of the following is/are present in mixture of product :
\qfig{p1_chem_s1_q04_1}
\begin{qoptions}
\opt{A}{$\mathrm{CH_3}$--$\mathrm{CH_2}$--OH}
\opt{B}{\qfigi{p1_chem_s1_q04_2}}
\opt{C}{Ph--${}^{18}$OH}
\opt{D}{\qfigi{p1_chem_s1_q04_D}}
\end{qoptions}

\qn{5} Which of the following compound can show positive carbyl amine test
\begin{qoptions}
\opt{A}{\qfigi{p1_chem_s1_q05_A}}
\opt{B}{\qfigi{p1_chem_s1_q05_B}}
\opt{C}{$\mathrm{CH_3CH_2CH_2NH_2}$ (1-aminopropane, drawn as a zig-zag chain with terminal $\mathrm{NH_2}$)}
\opt{D}{\qfigi{p1_chem_s1_q05_D}}
\end{qoptions}

\qn{6} $\mathrm{Ph-\overset{O}{\overset{\|}{C}}-CH_3\ \xrightarrow{\ A\ }\ Ph\text{--}CH_2\text{--}CH_3}$\\
A could be:
\begin{qoptions}
\opt{A}{$\mathrm{NH_2NH_2}$, glycol/$\mathrm{OH^-}$, $\Delta$}
\opt{B}{Zn(Hg)/conc. HCl}
\opt{C}{Red P/HI}
\opt{D}{$\begin{smallmatrix}\mathrm{CH_2-CH_2}\\ |\quad\ \ |\\ \mathrm{SH\ \ SH}\end{smallmatrix}$ ; Raney Ni, $\mathrm{H_2}$, $\Delta$}
\end{qoptions}

\qn{7} Which of the following is/are correct statement for $\mathrm{[Cl_4PtCl_2Pt(en)_2]Cl_2}$ ?
\begin{qoptions}
\opt{A}{It can show geometrical isomerism}
\opt{B}{It can show co-ordination position isomerism}
\opt{C}{It can show optical activity}
\opt{D}{It can show polymerisation isomerism}
\end{qoptions}

\qn{8} Which of following order is correct.
\begin{qoptions}
\opt{A}{$\mathrm{K_2O_3}$ is a mixed oxide}
\opt{B}{$\mathrm{K_2O}$ is used in submarine to absorb $\mathrm{CO_2}$}
\opt{C}{KOH is preferred to absorb $\mathrm{CO_2}$ compared to NaOH}
\opt{D}{$|\Delta H_f|$ order for alkali metals fluoride increases down the group}
\end{qoptions}

\sect{SECTION-II}
\qn{1} Pressure of 10 moles of an ideal gas is changed from 2 atm to 1 atm against constant external pressure without change in temperature. If surrounding temperature (300 K) and pressure (1 atm) always remains constant then calculate total entropy change ($\Delta S_{\text{system}}+\Delta S_{\text{surrounding}}$) for given process (J/K).\\
{[Given : $\ln 2=0.70$ and $R=8.0$ J/mol/K]}

\qn{2} Decomposition of A follows first order kinetics by the following equation.\\
$4\mathrm{A(g)}\to \mathrm{B(g)}+2\mathrm{C(g)}$\\
If initially, total pressure was 800 mm of Hg and after 10 minutes total pressure is found to be 650 mm of Hg. What is half-life of A (in minutes)? (Assume only A is present initially)

\qn{3} Adsorption of hydrogen gas over palladium surface obey Freundlich's adsorption isotherm as shown in the graph. The mass (in g) of hydrogen gas adsorbed on the surface of 1.0 g palladium at a pressure of $8.1\times10^{-3}$ atm is
\qfig{p1_chem_s2_q03_1}

\qn{4} A structure of hexapeptide can have how many peptide linkage :

\qn{5} NaCl aq. solution $\xrightarrow[\text{electrolysis}]{\text{(Inert Electrode)}}$ Gas-1 (at anode) + Gas-2 (cathode) + solution (S1)\\
Gas-1 + KI$_{\text{aq.}}$ $\longrightarrow$ diatomic molecule (X) + solution (S2)\\
Solution (S1) + salt (Y) $\longrightarrow$ Gas-3 + solution (S-3)\\
Salt-(Y) + conc. $\mathrm{H_2SO_4}$ $\longrightarrow$ Gas-4 + other products\\
Gas-3 + Gas-4 $\longrightarrow$ white fumes of Salt-Y\\
Find sum of the molecular weight of \textbf{Gas-1 + Gas-4 + diatomic compound (X) + Salt(Y)}\\
Given atomic weight of H = 1, Cl = 35.5, N = 14, I = 127, Na = 23, K = 39, S = 32, O = 16

\qn{6} Find sum of the number of atoms in the specified products of each molecule in given reactions.\\
(a) $\mathrm{K_{(s)}}+\underset{\text{Excess}}{\mathrm{O_2}}\longrightarrow$ Principal combustion product\\
(b) $\mathrm{KCl_{(Aq.)}}$ + hexachloridoplatinic(IV) acid $\longrightarrow$ Precipitate formed in the reaction\\
(c) $\mathrm{MnSO_4(aq.)}$ + sodium peroxydisulphate $\xrightarrow{\text{acid}}$ ``Mn'' containing ion\\
(d) $\mathrm{Na_2O_{2(s)}}+\mathrm{CO_{2(g)}}\longrightarrow$ gaseous product of reaction

\qn{7} Consider MOT for $\mathrm{B_2}$ \& $\mathrm{F_2}$ molecules \& find total number of electrons in both which is/are present in M.O. which is/are symmetrical around bond axis :-

\qn{8} Find number of process or ore given in Column-I are correctly represented with Column-II. You have to consider P for P, Q for Q, R for R etc.
\begin{center}\footnotesize
\begin{tabular}{|p{0.65cm}|p{2.5cm}|p{0.65cm}|p{2.5cm}|}
\hline
\multicolumn{2}{|c|}{(I) Process/Ore} & \multicolumn{2}{c|}{(II) Related to}\\ \hline
(P) & Matte formation & (P) & Cu\\ \hline
(Q) & Cynide leaching & (Q) & Au\\ \hline
(R) & Monds process & (R) & Reaction of metal + $\mathrm{H_2}$\\ \hline
(S) & Fused oxide electrolysis & (S) & Al\\ \hline
(T) & Chromite ore & (T) & Fe + Cr\\ \hline
(U) & Van Arkel process & (U) & Zr\\ \hline
(V) & Bacterial leaching & (V) & Low grade Cu ore\\ \hline
(W) & Separation of $\mathrm{FeWO_4}$ & (W) & Extraction of Pb\\ \hline
\end{tabular}
\end{center}

\sect{SECTION-IV}
\begin{center}\footnotesize
\begin{tabular}{|p{0.6cm}|p{6.6cm}|}
\hline
\multicolumn{2}{|c|}{Column - I (Chemical reaction)}\\ \hline
(A) & \qfigi{p1_chem_s3_pre_A}\\ \hline
(B) & \qfigi{p1_chem_s3_pre_B}\\ \hline
(C) & \qfigi{p1_chem_s3_pre_C}\\ \hline
(D) & \qfigi{p1_chem_s3_pre_D}\\ \hline
\end{tabular}
\end{center}
\qn{1} 
\begin{center}\footnotesize
\begin{tabular}{|p{0.6cm}|p{6.6cm}|}
\hline
\multicolumn{2}{|c|}{Column - II (Type of reaction/intermediate involve)}\\ \hline
(P) & Elimination reaction\\ \hline
(Q) & Addition reaction\\ \hline
(R) & Carbocation intermediate is formed [Classical /Non classical]\\ \hline
(S) & Rearrangment reaction\\ \hline
(T) & Racemic mixture is formed\\ \hline
\end{tabular}
\end{center}

\qn{2} 
\begin{center}\footnotesize
\begin{tabular}{|p{0.55cm}|p{3.2cm}|p{0.6cm}|p{2.9cm}|}
\hline
\multicolumn{2}{|c|}{Column - I} & \multicolumn{2}{c|}{Column - II}\\ \hline
(A) & Ethane & (P) & Gauche form is more stable than anti form\\ \hline
(B) & Propane & (Q) & Anti form is more stable than gauche form\\ \hline
(C) & 1,2-Dichloroethane & (R) & Can show conformational isomerism\\ \hline
(D) & 1-Chloropropane & (S) & Can show position isomerism\\ \hline
 & & (T) & Only one stable conformation\\ \hline
\end{tabular}
\end{center}
\partsec{3}{MATHEMATICS}
\sect{SECTION-I}
\qn{1} In three dimensional space, the equation
\[\resizebox{\linewidth}{!}{$\begin{vmatrix}1+2x+3x^2&2+3x+x^2&3+x+2x^2\\1+2y+3y^2&2+3y+y^2&3+y+2y^2\\1+2z+3z^2&2+3z+z^2&3+z+2z^2\end{vmatrix}=0$}\]
represents
\begin{qoptions}
\opt{A}{Three planes}
\opt{B}{Only three straight lines}
\opt{C}{A sphere}
\opt{D}{Only a point}
\end{qoptions}

\qn{2} Let $M$ be the locus of the point lying in the plane $P: x+y+z=3$, whose distance from origin is always $\sqrt5$ unit. The curve $M$ meets the plane $y=0$ at two points $A$ and $B$ (abscissa of $A$ is greater than that of $B$). Line $L$ lying in the plane $P$ meets the curve $M$ at only one point $A$. Then choose the correct statement(s)
\begin{qoptions}
\opt{A}{Coordinates of point $A$ is $(1,0,2)$}
\opt{B}{Length of the curve $M$ is $2\sqrt2\,\pi$}
\opt{C}{Direction ratio of line $L$ is $(1,1,-2)$}
\opt{D}{Distance of line $L$ from origin is $\sqrt5$ unit}
\end{qoptions}

\qn{3} Let $y(x)$ be a solution of the differential equation $\left(x-\dfrac{dy}{dx}\right)\sin x=y\cos x,\ x\in\left(0,\dfrac\pi2\right]$. If $y\!\left(\dfrac\pi2\right)=1$, then choose the correct statement(s) :-
\begin{qoptions}
\opt{A}{$y(x)$ has exactly one maxima in $\left(0,\dfrac\pi2\right]$}
\opt{B}{$y(x)$ has exactly one minima in $\left(0,\dfrac\pi2\right]$}
\opt{C}{$\displaystyle\lim_{x\to0^+}y(x)=0$}
\opt{D}{$0<y\!\left(\dfrac\pi4\right)<1$}
\end{qoptions}

\qn{4} Let $f(x)=\begin{cases}x\ln|x|; & x\ne0\\ 0; & x=0\end{cases}$ and $g(x)=x\{x\}+1;\ x\ge-1$. Then, choose correct statement(s) :- (where $\{.\}$ is fractional part)
\begin{qoptions}
\opt{A}{$f(x)$ is differentiable at $x=0$}
\opt{B}{$g(x)$ is differentiable at $x=0$}
\opt{C}{$f\circ g(x)$ is not differentiable at $x=0$}
\opt{D}{$f\circ g(x)$ is not differentiable at $x=1$}
\end{qoptions}

\qn{5} If $\alpha=4\tan^{-1}(\sqrt3-1)$ and $\beta=8\cot^{-1}2$, then choose the correct statement(s) :-
\begin{qoptions}
\opt{A}{$\cos\alpha<0$}
\opt{B}{$\tan\alpha>0$}
\opt{C}{$0<\tan\beta<\sqrt3$}
\opt{D}{$\tan^{-1}\dfrac{47}{48}<\tan^{-1}\dfrac{48}{49}$}
\end{qoptions}

\qn{6} Let curve $C_1$ be the locus of mid-points of all chords of ellipse $E: x^2+2y^2=2$, that passes through the point $Q(-2,0)$. Further, curve $C_2$ is the locus of the point $P$ such that $QA$, $QP$ and $QB$ are in A.P., where the line $QP$ meets the ellipse $E$ at points $A$ and $B$. Now, choose the correct statement(s) :-
\begin{qoptions}
\opt{A}{Curve $C_1$ is a part of ellipse $x^2+2x+2y^2=0$}
\opt{B}{Curve $C_2$ is a part of ellipse $x^2+2x+2y^2=0$}
\opt{C}{Putting aside origin, no point of curve $C_1$ lies in I$^{\text{st}}$ quadrant}
\opt{D}{Putting aside origin, no point of curve $C_2$ lies in IV$^{\text{th}}$ quadrant}
\end{qoptions}

\qn{7} Let $A$ and $B$ be two respectively symmetric and skew symmetric, $3\times3$, non-zero, matrices and $P$ be an arbitrary $3\times3$, non-zero matrix. Then which of the following is/are symmetric matrix ?
\begin{qoptions}
\opt{A}{$P^TP+B^2$}
\opt{B}{$P^TAP+A$}
\opt{C}{$(AB+BA)^2$}
\opt{D}{$(P^TBP)^2+AP+P^TA$}
\end{qoptions}

\qn{8} Let $f(x)=e^{\sin x}-e^{-\sin x};\ x\in\left[0,\dfrac\pi2\right]$ and $g(x)=\dfrac\pi2(1-\cos x);\ x\in\left[0,\dfrac\pi2\right]$. Then, which of the following is/are true ?
\begin{qoptions}
\opt{A}{Range of $f$ is $[0,e]$}
\opt{B}{Range of $f(g(x))$ is $\left[0,\ e-\dfrac1e\right]$}
\opt{C}{$\displaystyle\lim_{x\to0^+}\dfrac{f(x)}{g(x)}=\dfrac4\pi$}
\opt{D}{There exists atleast one real root of equation $g(f(x))=\dfrac\pi2$, in $x\in\left[0,\dfrac\pi2\right]$}
\end{qoptions}

\sect{SECTION-II}
\qn{1} The radius of the largest circle whose center lies on the $y$-axis and which passes through the focus of parabola $x^2=4y$ and touching it, is

\qn{2} A multibarrel gun fires 7 bullets in one round. A self-repairable target is destroyed when it is hit by atleast two bullets in any single round. Further, a bullet hits the target with probability $\dfrac12$. Then, the minimum number of continuous rounds must be fired so that the probability of destroying the target is atleast $\dfrac{255}{256}$, is

\qn{3} Let $f:\mathbb R\to\mathbb R$, be a function defined by $f(x)=\begin{cases}[x]; & x<0\\ x; & x\ge0\end{cases}$; where $[x]$ is greatest integer function. Also, $g(x)=(x+1)(x+2)(x+3)$. Now, the value of $\displaystyle\int_{-3}^{2}g\big(f(x)\big)\,dx$, is

\qn{4} A closed rectangular box with square base has been made from certain solid material. Its inner volume is 512 cubic unit and all six walls are 2 unit thick. Now, some rectangular openings through the wall is carved out, whose area is half of the initial inner surface area of box. Then, the material volume of this partially-closed box is minimum, when height of the box is

\qn{5} The number of distinct real solutions of equation $4\sin^4x+4\sin^3x+3\sin^2x=1+\sin x$, in the interval $[0,2\pi]$, is

\qn{6} The value of $\dfrac{64}{\pi^2}\displaystyle\int_0^1\dfrac{\tan^{-1}\!\left(\frac{2x}{x+1}\right)}{5x^2+2x+1}\,dx$

\qn{7} If $\alpha_k$; $k=1,2,\ldots,6$; are six distinct sixth roots of unity, then the value of $\left(\displaystyle\sum_{1\le i<j\le6}|\alpha_j-\alpha_i|\right)(2-\sqrt3)$ is equal to

\qn{8} If coefficient of $x^2$ in the expansion of $(1+2x+x^2+2x^3+x^4+2x^5+\ldots+2x^{2019})^{50}(1-x+x^3-x^5+x^7-\ldots+x^{2019})^{100}$ is $k$, then the value of $(k+200)$ is equal to

\sect{SECTION-IV}
\qn{1} 
\begin{center}\footnotesize
\begin{tabular}{|p{0.5cm}|p{5.0cm}|p{0.6cm}|p{0.7cm}|}
\hline
\multicolumn{2}{|c|}{Column-I} & \multicolumn{2}{c|}{Column-II}\\ \hline
(A) & If projections of a unit vector $\hat r$ along the vectors $\dfrac{\hat i+\hat j}{\sqrt2},\ \dfrac{-\hat i+\hat j}{\sqrt2}$ and $\hat k$ are equal. Then, the value of $\sqrt3\,\hat r\cdot(\hat i+\hat k)$ can be & (P) & $-1$\\ \hline
(B) & The number of points at which function $f:\mathbb R\to\mathbb R$, $f(x)=\sqrt[3]{x^2|x|}-|x|-1$, is not differentiable, is & (Q) & 0\\ \hline
(C) & The remainder when the sum $\displaystyle\sum_{r=1}^{10}(r+1)(r\omega+1)(r\omega^2+1)$, is divided by 3, is greater than (where $\omega$ is imaginary cube root of unity) & (R) & 1\\ \hline
(D) & The cubes of the natural number are grouped as $1^3,\ (2^3,3^3),\ (4^3,5^3,6^3)\ldots$. Let $S_n$ denotes the sum of cubes in $n$th group, then sum of the digits of $8S_{10}$ is greater than & (S) & 2\\ \hline
 & & (T) & 5\\ \hline
\end{tabular}
\end{center}

\qn{2} 
\begin{center}\footnotesize
\begin{tabular}{|p{0.5cm}|p{5.0cm}|p{0.6cm}|p{0.7cm}|}
\hline
\multicolumn{2}{|c|}{Column-I} & \multicolumn{2}{c|}{Column-II}\\ \hline
(A) & The area of the triangle $ABC$ is $\sqrt3$ sq. unit and $\angle A=\dfrac{2\pi}3$, then the maximum value of $(b\sin C+c\sin B)$ is & (P) & $\dfrac12$\\ \hline
(B) & In $\triangle ABC$, it is given that $b=3c$, then the value of $\cot\!\left(\dfrac{B-C}2\right)\cot\!\left(\dfrac A2\right)$ is & (Q) & $\dfrac34$\\ \hline
(C) & The volume of parallelopiped whose three coterminous edges are $\sin\theta\,\hat i-\sin\theta\,\hat k$, $2\hat i+2\cos\theta\,\hat j$ and $\hat i+\sin\theta\,\hat j+\hat k$, can be & (R) & 1\\ \hline
(D) & The area of the smaller region bounded by the circle $x^2+y^2=1$ and the parabola $(x-y)^2=x+y$, is less than or equal to & (S) & 2\\ \hline
 & & (T) & 3\\ \hline
\end{tabular}
\end{center}
\end{multicols}
\end{document}
